
\documentclass{article}
\usepackage{colortbl}
\usepackage{makecell}
\usepackage{multirow}
\usepackage{supertabular}

\begin{document}

\newcounter{utterance}

\centering \large Interaction Transcript for game `cladder', experiment `full\_v1.5\_default', episode 7046 with llama{-}3.1{-}8b{-}instruct.
\vspace{24pt}

{ \footnotesize  \setcounter{utterance}{1}
\setlength{\tabcolsep}{0pt}
\begin{supertabular}{c@{$\;$}|p{.15\linewidth}@{}p{.15\linewidth}p{.15\linewidth}p{.15\linewidth}p{.15\linewidth}p{.15\linewidth}}
    \# & \multicolumn{2}{c}{Player} && \multicolumn{2}{c}{Game Master} \\
    \hline

    \theutterance \stepcounter{utterance}  
    & & & \multicolumn{4}{p{0.6\linewidth}}{
        \cellcolor[rgb]{0.9,0.9,0.9}{
            \makecell[{{p{\linewidth}}}]{
                \texttt{\tiny{[P1$\langle$GM]}}
                \texttt{You are an expert in causal inference. The following question is not a typical commonsense query, but rather a meticulously designed question created by a professor specializing in causal inference, intended to assess the students' mastery of the course content.} \\
\texttt{Imagine a self{-}contained, hypothetical world with only the following conditions, and without any unmentioned factors or causal relationships: Jyka has a direct effect on lirg. Hwax has a direct effect on lirg. The overall probability of jyka is 14\%. For those who are not jyka, the probability of lirg is 71\%. For those who are jyka, the probability of lirg is 66\%. Is lirg less likely than not lirg overall?} \\
\texttt{Start your answer with 'Yes' or 'No', followed by additional reasoning or evidence to support your explanation.} \\
            }
        }
    }
     \\ \\

    \theutterance \stepcounter{utterance}  
    & \multicolumn{4}{p{0.6\linewidth}}{
        \cellcolor[rgb]{0.561,0.737,0.561}{
            \makecell[{{p{\linewidth}}}]{
                \texttt{\tiny{[P1$\rangle$GM]}}
                \texttt{No.} \\
\\ 
\texttt{To determine whether lirg is less likely than not lirg overall, we need to calculate the overall probability of lirg and compare it to the probability of not lirg.} \\
\\ 
\texttt{Let's denote the probability of jyka as P(J) = 0.14, and the probability of lirg given jyka as P(L|J) = 0.66, and the probability of lirg given not jyka as P(L|\textasciitilde{}J) = 0.71.} \\
\\ 
\texttt{We can use the law of total probability to calculate the overall probability of lirg:} \\
\\ 
\texttt{P(L) = P(L|J) * P(J) + P(L|\textasciitilde{}J) * P(\textasciitilde{}J)} \\
\texttt{= 0.66 * 0.14 + 0.71 * (1 {-} 0.14)} \\
\texttt{= 0.0924 + 0.6136} \\
\texttt{= 0.706} \\
\\ 
\texttt{Now, we need to calculate the probability of not lirg:} \\
\\ 
\texttt{P(\textasciitilde{}L) = 1 {-} P(L)} \\
\texttt{= 1 {-} 0.706} \\
\texttt{= 0.294} \\
\\ 
\texttt{Since 0.706 is greater than 0.294, lirg is more likely than not lirg overall.} \\
            }
        }
    }
    & & \\ \\

    \theutterance \stepcounter{utterance}  
    & & & \multicolumn{2}{p{0.3\linewidth}}{
        \cellcolor[rgb]{0.9,0.9,0.9}{
            \makecell[{{p{\linewidth}}}]{
                \texttt{\tiny{[GM$|$GM]}}
                \texttt{No} \\
            }
        }
    }
    & & \\ \\

    \theutterance \stepcounter{utterance}  
    & & & \multicolumn{2}{p{0.3\linewidth}}{
        \cellcolor[rgb]{0.9,0.9,0.9}{
            \makecell[{{p{\linewidth}}}]{
                \texttt{\tiny{[GM$|$GM]}}
                \texttt{no} \\
            }
        }
    }
    & & \\ \\

    \theutterance \stepcounter{utterance}  
    & & & \multicolumn{2}{p{0.3\linewidth}}{
        \cellcolor[rgb]{0.9,0.9,0.9}{
            \makecell[{{p{\linewidth}}}]{
                \texttt{\tiny{[GM$|$GM]}}
                \texttt{game\_result = WIN} \\
            }
        }
    }
    & & \\ \\

\end{supertabular}
}

\end{document}
